% Part: first-order-logic
% Chapter: models-theories
% Section: size-of-structures

\documentclass[../../../include/open-logic-section]{subfiles}

\begin{document}

\olfileid{fol}{mat}{siz}

\olsection{بیان اندازهٔ \printtoken{P}{structure}}

\begin{explain}
برخی خواص ساختارها را حتی بدون استفاده از نمادهای غیرمنطقی یک زبان
می‌توانیم بیان کنیم. برای نمونه، !!{sentence}sی وجود دارند که در
!!a{structure} صادق‌اند هرگاه و تنها هرگاه !!{domain} آن !!{structure}
دست‌کم، حداکثر، یا دقیقاً تعداد معینی چون~$n$ !!{element} داشته باشد.
\end{explain}

\begin{prop}
!!{sentence}
\begin{multline*}
  !A_{\ge n} \ident \lexists[x_1][\lexists[x_2][\dots\lexists[x_n][{}]]]\\
  \begin{aligned}
  (\eq/[x_1][x_2] \land {}
  \eq/[x_1][x_3] \land \eq/[x_1][x_4] \land \dots \land \eq/[x_1][x_n] \land {}\\
\eq/[x_2][x_3] \land \eq/[x_2][x_4] \land \dots \land {} \eq/[x_2][x_n] \land {} \\
\vdots\\
\eq/[x_{n-1}][x_n])
\end{aligned}
\end{multline*}
در !!a{structure}~$\Struct M$ صادق است هرگاه و تنها هرگاه~$\Domain M$
دست‌کم $n$~!!{element}s داشته باشد. در نتیجه،
$\Sat{M}{\lnot !A_{\ge n+1}}$ هرگاه و تنها هرگاه~$\Domain M$ حداکثر
$n$~!!{element}s داشته باشد.

\end{prop}

\begin{prop}
!!{sentence}
\begin{multline*}
!A_{= n} \ident \lexists[x_1][\lexists[x_2][\dots\lexists[x_n][{}]]] \\
\begin{aligned}
  (\eq/[x_1][x_2] \land {}
  \eq/[x_1][x_3] \land \eq/[x_1][x_4] \land \dots \land \eq/[x_1][x_n] \land {}\\
\eq/[x_2][x_3] \land \eq/[x_2][x_4] \land \dots \land {} \eq/[x_2][x_n] \land {} \\
\vdots\\
\eq/[x_{n-1}][x_n] \land {} \\
\lforall[y][(\eq[y][x_1] \lor \dots \lor \eq[y][x_n])]
\end{aligned}
\end{multline*}
در !!a{structure}~$\Struct M$ صادق است هرگاه و تنها هرگاه~$\Domain M$
دقیقاً $n$~!!{element}s داشته باشد.
\end{prop}

\begin{prop}
!!a{structure} نامتناهی است هرگاه و تنها هرگاه مدلی از مجموعهٔ زیر باشد:
\[
\{!A_{\ge 1}, !A_{\ge 2}, !A_{\ge 3}, \dots \}.
\]
\end{prop}

هیچ جملهٔ صرفاً منطقیِ واحدی وجود ندارد که در~$\Struct M$ صادق
باشد هرگاه و تنها هرگاه~$\Domain M$ نامتناهی باشد. بااین‌حال، می‌توان
!!{sentence}sی دارای !!{predicate}s غیرمنطقی به دست داد که فقط مدل‌های
نامتناهی دارند (هرچند هر !!{structure} نامتناهی‌ای مدل آن‌ها نیست).
خاصیت متناهی‌بودن یک ساختار و خاصیت !!{nonenumerable}بودن یک ساختار را
حتی با مجموعه‌ای نامتناهی از !!{sentence}s نیز نمی‌توان بیان کرد. این
واقعیت‌ها از قضیه‌های فشردگی و لوونهایم--اسکولم نتیجه می‌شوند.

\end{document}
